\section{Introduction}
\label{sec:Introduction}

Dans la musique contemporaine, certaines partitions sont écrites par des 
personnes qui ne pratiquent pas l'instrument. Il peut arriver que certains 
compositeurs demandent à des musiciens de jouer des partitions difficiles, 
voire impossibles. Aujourd'hui, le problème peut aussi se poser avec l'arrivé
d'outils de génération probabiliste de partitions, ou de suite de positions.
C'est sur cette considération que Charlotte \textsc{Truchet} et Mikhail 
\textsc{Malt} ont décidés d'étudier la jouabilité de partition de guitare dans
le cadre de la thèse de Charlotte \cite{CTThese}. Le problème était le 
suivant : étant donné une suite d'accords, est-il possible de l'interpréter à 
la guitare, si oui, comment l'interpréter de la manière la plus facile ?\\
Au cours de mon stage, j'ai repris cette question dans l'objectif de réaliser
un programme permettant de calculer une suite de positions à partir d'une suite
d'accords en utilisant une approche par programation sous contraintes. 
Pour l'implémentation de la résolutions, j'ai utilisé le solver Choco
développéé par Charles \textsc{Prud'homme} \cite{Choco}.
Dans un premier temps, je me suis familiarisé avec la notion de programmation
sous contraintes, et ensuite, j'ai réalisé un modèle pour résoudre le problème
posé, en discutant avec des guitaristes (ne l'étant pas moi même) afin de
laisser le plus de liberté et de maniabilité aux instrumentistes.
